\section{Краевые задачи для дифференциальных уравнений в частных производных}
\label{sec:boundary-value-problems-for-pdes}

Уравнения в частных производных описывают поля, зависящие от нескольких
пространственных переменных и, возможно, времени. Для однозначного
определения решения одного дифференциального уравнения недостаточно:
необходимо задать начальные и/или граничные условия.

\subsection{Классификация уравнений второго порядка}

Для линейного уравнения двух переменных
\[
A u_{xx}+2B u_{xy}+C u_{yy}
+\text{члены низшего порядка}=f
\]
тип в точке определяется знаком
\[
D=B^2-AC.
\]
\begin{itemize}
    \item $D<0$ --- \textbf{эллиптический тип}; пример: уравнение Лапласа
    или Пуассона.
    \item $D=0$ --- \textbf{параболический тип}; типичный пример ---
    уравнение теплопроводности.
    \item $D>0$ --- \textbf{гиперболический тип}; типичный пример ---
    волновое уравнение.
\end{itemize}
Для уравнений более высокой размерности классификация определяется
сигнатурой квадратичной формы главной части.

\subsection{Начальные и граничные условия}

На границе $\partial\Omega$ часто используются:
\begin{itemize}
    \item условие Дирихле
    \[
    u=g;
    \]
    \item условие Неймана
    \[
    \frac{\partial u}{\partial n}=g;
    \]
    \item условие Робена
    \[
    \alpha u+\beta\frac{\partial u}{\partial n}=g.
    \]
\end{itemize}
Для нестационарных задач дополнительно задаются начальные условия. Их
число связано с порядком производной по времени.

\subsection{Эллиптические задачи: Лаплас и Пуассон}

Модельная задача Дирихле:
\[
-\Delta u=f\quad\text{в }\Omega,
\qquad
u=0\quad\text{на }\partial\Omega.
\]
Для гармонической функции ($\Delta u=0$) действует \textbf{принцип
максимума}: максимум и минимум непрерывного решения на ограниченной
области достигаются на границе, если решение не является постоянным.
В частности, из него следует единственность задачи Дирихле.

Для задачи Неймана
\[
-\Delta u=f,
\qquad
\frac{\partial u}{\partial n}=g
\]
необходимо условие совместимости
\[
\boxed{
\int_\Omega f\,dx+\int_{\partial\Omega}g\,ds=0.
}
\]
При его выполнении решение определяется с точностью до добавления
постоянной.

\subsection{Слабая постановка и теорема Лакса--Мильграма}

Для однородной задачи Дирихле естественным пространством является
$V=H_0^1(\Omega)$. После умножения на тестовую функцию $v$ и интегрирования
по частям получаем слабую постановку:
\[
\boxed{
\text{найти }u\in V:\qquad
\int_\Omega \nabla u\cdot\nabla v\,dx
=\int_\Omega fv\,dx
\quad\forall v\in V.
}
\]

Общий результат задаётся \textbf{теоремой Лакса--Мильграма}. Если $V$ ---
гильбертово пространство, билинейная форма $a$ непрерывна и коэрцитивна,
\[
|a(u,v)|\le M\|u\|_V\|v\|_V,
\qquad
a(v,v)\ge\alpha\|v\|_V^2,
\quad \alpha>0,
\]
а $\ell\in V^*$, то существует единственное $u\in V$, удовлетворяющее
\[
a(u,v)=\ell(v)\qquad\forall v\in V.
\]
Для задачи Пуассона коэрцитивность следует из неравенства Пуанкаре.
Эта теорема даёт строгую основу существования и единственности слабого
решения и лежит в основе МКЭ.

\subsection{Параболические задачи}

Уравнение теплопроводности имеет вид
\[
\frac{\partial u}{\partial t}-a^2\Delta u=f,
\qquad a>0.
\]
Для него обычно задают одно начальное условие
\[
u(x,0)=u_0(x)
\]
и граничные условия на боковой границе. Параболические уравнения обладают
сглаживающим свойством: локальные возмущения распространяются по области и
со временем высокочастотные компоненты затухают.

Аналог принципа максимума даёт оценки решения через начальные, граничные
данные и правую часть и используется для доказательства единственности и
устойчивости.

\subsection{Гиперболические задачи}

Для волнового уравнения
\[
u_{tt}-c^2\Delta u=f
\]
нужно задать два начальных условия
\[
u(x,0)=u_0(x),
\qquad
u_t(x,0)=v_0(x),
\]
а также соответствующие граничные условия.

Для однородного волнового уравнения с однородными граничными условиями
энергия
\[
\boxed{
E(t)=\frac12\int_\Omega
\left(u_t^2+c^2|\nabla u|^2\right)dx
}
\]
сохраняется. Энергетическая оценка является одним из основных средств
доказательства единственности и непрерывной зависимости решения от данных.
В отличие от параболического случая возмущения распространяются с конечной
скоростью, связанной с характеристиками.

\subsection{Корректность по Адамару}

Задача называется \textbf{корректно поставленной}, если
\begin{enumerate}
    \item решение существует;
    \item решение единственно;
    \item решение непрерывно зависит от исходных данных.
\end{enumerate}
Нарушение любого из этих требований делает задачу некорректной. Перед
численным решением важно понимать, корректна ли сама непрерывная
постановка: вычислительный метод не может устранить фундаментальную
неустойчивость исходной задачи без дополнительной регуляризации.

\subsection{Численные методы}

Основные подходы:
\begin{itemize}
    \item метод конечных разностей;
    \item метод конечных элементов;
    \item метод конечных объёмов;
    \item спектральные методы;
    \item метод характеристик для гиперболических задач.
\end{itemize}
Выбор метода определяется типом уравнения, геометрией области,
регулярностью решения и требованием сохранения физических законов.


\subsection{Формулы Грина и происхождение слабой постановки}

Для достаточно гладких $u$ и $v$ первая формула Грина имеет вид
\[
    \int_\Omega \nabla u\cdot\nabla v\,dx
    =-\int_\Omega v\,\Delta u\,dx
     +\int_{\partial\Omega}v\frac{\partial u}{\partial n}\,ds.
\]
Именно она переносит одну производную с решения на тестовую функцию и позволяет
снизить требования к гладкости решения. Условие Дирихле включается в пространство
допустимых функций, а условие Неймана возникает в граничном интеграле как
естественное условие.

Для более общего эллиптического оператора
\[
    -\nabla\cdot(A(x)\nabla u)+c(x)u=f
\]
при симметричной матрице $A$ и условии равномерной эллиптичности
\[
    \xi^T A(x)\xi\ge \alpha|\xi|^2
\]
слабая форма имеет вид
\[
    a(u,v)=\ell(v),
\]
\[
    a(u,v)=\int_\Omega
    \left(A\nabla u\cdot\nabla v+cuv\right)dx.
\]
При соответствующих условиях на $c$ и граничную часть Дирихле эта форма
коэрцитивна, и теорема Лакса--Мильграма даёт существование и единственность.

\subsection{Неоднородное условие Дирихле и подъём граничных данных}

Если
\[
    u=g\quad\text{на }\Gamma_D,
\]
то пространство допустимых функций является аффинным. Выбирают продолжение
$u_g\in H^1(\Omega)$ с нужным следом и представляют
\[
    u=u_g+w,
    \qquad w\in H^1_0(\Omega;\Gamma_D).
\]
После подстановки задача сводится к слабой задаче относительно $w$ с однородным
условием Дирихле. Этот простой приём важен и в теории, и в МКЭ.

\subsection{Энергетическая оценка для теплопроводности}

Для однородной задачи
\[
    u_t-a^2\Delta u=f,
    \qquad u|_{\partial\Omega}=0
\]
умножение на $u$ и интегрирование дают
\[
    \frac12\frac{d}{dt}\|u\|_{L^2}^2
    +a^2\|\nabla u\|_{L^2}^2
    =(f,u).
\]
При $f=0$ норма $L^2$ не возрастает, а градиент диссипирует энергию. Такая
оценка объясняет устойчивость и сглаживающее действие параболического уравнения.
С помощью неравенств Коши и Гронуолла получают непрерывную зависимость от данных.

\subsection{Характеристики и конечная скорость распространения}

Для гиперболических уравнений информация распространяется вдоль
характеристик. В одномерном волновом уравнении
\[
    u_{tt}-c^2u_{xx}=0
\]
характеристики имеют вид
\[
    x\pm ct=\text{const}.
\]
Следовательно, значение решения в точке зависит только от данных внутри
соответствующего конуса зависимости. Это принципиально отличает гиперболические
задачи от уравнения теплопроводности, где математическое влияние мгновенно
распространяется по всей области.

\subsection{Связь корректности и численного метода}

Численная схема должна наследовать ключевые свойства непрерывной задачи.
Для эллиптических задач важны дискретная коэрцитивность или принцип максимума,
для параболических и гиперболических --- устойчивость по времени и корректная
обработка начальных и граничных условий. Метод конечных разностей напрямую
аппроксимирует сильную форму, а метод конечных элементов обычно строится из
слабой формы. Поэтому выбор численного метода тесно связан с типом уравнения и
естественной функциональной постановкой задачи.

\subsection{Что важно сказать на экзамене}

Нужно назвать три типа УЧП второго порядка, привести по одному модельному
уравнению и объяснить, какие начальные/граничные данные требуются. Для
эллиптической задачи полезно сформулировать принцип максимума и слабую
постановку с теоремой Лакса--Мильграма; для волновой задачи --- энергетическую
оценку. Обязательно знать определение корректности по Адамару.

\newpage
